%% ma: L12-B01-0024
%% lop: 12
%% chuong: 1
%% bai: 1
%% dang: 12.1.16
%% loai: TN4
%% muc: TH
%% dap_an: "D"
%% nguon: "(NBV)-ĐỀ SỐ 13-MỖI NGÀY 1 ĐỀ THI 2025.pdf (D13, MỖI NGÀY 1 ĐỀ - Đề số 13), Phần I Câu 6"
%% kien_thuc: "Bài 1 (tính đơn điệu của hàm số qua đạo hàm)"
%% trang_thai: chinh-thuc
%% ghi_chu: "Đã duyệt 14/10/2026. Đáp án gốc D đúng (y' = 7/(x+2)^2 dương trên từng khoảng của tập xác định; hàm đồng biến trên mỗi khoảng (-vô cùng; -2) và (-2; +vô cùng)). Phương án B và C chứa điểm -2 hoặc hai khoảng rời nhau nên sai."
\begin{ex}
Hàm số $y=\dfrac{x-5}{x+2}$ đồng biến trên
\choice{$\mathbb R$}{$\mathbb R\setminus\{-2\}$}{$(-5;+\infty)$}{\True $(-\infty;-2)$}
\loigiai{Tập xác định $\mathbb R\setminus\{-2\}$. Ta có $y'=\dfrac{7}{(x+2)^2}>0$ với mọi $x\ne-2$, nên hàm số đồng biến trên mỗi khoảng $(-\infty;-2)$ và $(-2;+\infty)$.

Các tập $\mathbb R$ và $(-5;+\infty)$ chứa điểm $x=-2$ không thuộc tập xác định, còn $\mathbb R\setminus\{-2\}$ gồm hai khoảng rời nhau (hàm số không đồng biến trên cả hợp hai khoảng đó vì $f(-3)=8>f(0)=-2{,}5$ dù $-3<0$). Chọn D.}
\end{ex}
